\itemhead{p.84 练习 · 第1题（教材p.84）}{XJ-B1-C03-PR-P084-Y1009}
\label{XJ-B1-C03-PR-P084-Y1009}
{\footnotesize 湘教版教材p.84 / PDF第91页 / p.84 练习 / 第1题}
\field{原题}{判断下列函数的奇偶性：
\subq{1}$f(x)=x+\dfrac{x^3}5-\dfrac{x^5}{10}$；
\subq{2}$f(x)=3-\dfrac{x^2}3$；
\subq{3}$f(x)=\dfrac{3x^2}{\sqrt{3+x^4}}$；
\subq{4}$f(x)=1-2x+x^3$。
这几个函数的图象如图所示，你能在图中分别标出对应的函数吗？
\bankfig{C03-PR084}{0.46}}
\field{参考答案}{依次为奇、偶、偶、非奇非偶；对应图中红色虚线、蓝色抛物线、红色实线、绿色曲线。可用在0处的值及轴对称性辅助核对。}
{\small \field{关联知识}{F08 奇偶性、对称与函数性质应用（教材帮p.113、114、115 / PDF122、123、124）；F04 函数三种表示法、图象与读数（教材帮p.85、86、93、94 / PDF94、95、102、103）}
\field{方法与易错点}{检查定义域，再代入 $-x$；求另一侧解析式时同时写出 $-x$ 所在区间。 易错：只看公式不看定义域；用一对点证明奇偶；负半轴代换遗漏最外层负号。}
\field{答案依据}{答案册相关题组：PDF第23页（册内页22）；按当前题库原题逐项重算，不按旧版题号机械对套。}}
